In this calculation put $\ell=\ln X$, $s=1+X$, and
$\beta=X/(X+2)$. The letter $r$ denotes the positive square-root kernel. The functions used in the soft master are
\begin{align}
 k&=\sqrt{\frac{1+XU^2}{s}},\qquad d=1-\beta^2U^2,\\
 a(U)&=\operatorname{artanh}(\beta U),\qquad
 f(U)=\operatorname{artanh}(U/k),\qquad 0<U<1.
\end{align}
Here $d$ is a scalar function, unrelated to the spacetime dimension.
The symbol $a$ without an argument in the unit-circle identities
below denotes a complex number, not the function $a(U)$.

\paragraph{Root-of-unity endpoints.} Let $\theta=\pi/3$,
$a=\exp(\ii\theta)$, and define the real Clausen constant
\begin{equation}
 \mathcal C=\operatorname{Im}\operatorname{Li}_2(a)
           =\sum_{n=1}^{\infty}\frac{\sin(n\pi/3)}{n^2}.
\end{equation}
The rationalized root primitives have endpoints $q=1$ and $q=0$.
Their unit-circle values are obtained from
\begin{align}
 \operatorname{Li}_2(e^{\ii t})
 &=\frac{\pi^2}{6}-\frac{\pi|t|}{2}+\frac{t^2}{4}
       +\ii\operatorname{Cl}_2(t),\\
 \operatorname{Im}\operatorname{Li}_3(e^{\ii t})
 &=\frac{\pi^2t}{6}-\frac{\pi t|t|}{4}+\frac{t^3}{12},
                         \qquad -2\pi\leq t\leq2\pi.
\end{align}
Here $\operatorname{Cl}_2(t)=\sum_{n\geq1}\sin(nt)/n^2$ is the
Clausen function. These Fourier identities follow by integrating the geometric-series
logarithm, with constants fixed at $t=0$. Splitting the defining
series into residue classes also gives
\begin{equation}
 \operatorname{Re}\operatorname{Li}_3(a)=\frac{\zeta(3)}3,
 \quad \operatorname{Re}\operatorname{Li}_3(a^2)=-\frac{4\zeta(3)}9,
 \quad \operatorname{Cl}_2(2\theta)=\frac23\mathcal C.
\end{equation}
For explicit endpoint bookkeeping, write
\begin{equation}
 J_{mn}=(\ii m\theta)^2\ln(1-a^{m-n})
       +2\ii m\theta\operatorname{Li}_2(a^{m-n})
       -2\operatorname{Li}_3(a^{m-n}),\qquad J_{00}=-2\zeta(3).
\end{equation}
This is the earlier logarithm-squared primitive at a unit-circle
endpoint. All phases are the continuous ones inherited from $0<q<1$.
For $m=\pm1,\pm2$, the principal phases give those same endpoint
values. With $S_3=\operatorname{Li}_3(a)+\operatorname{Li}_3(a^{-1})$,
the constants of the corresponding primitives are
\begin{align}
 P_A&=2J_{-1,0}+2J_{1,0}-J_{-2,0}+J_{-2,2}-2S_3,\\
 M_n&=\frac12\left[-2\zeta(3)+(-\ii n\theta)^2(-2\ii n\theta)
                  -4\ii n\theta\zeta(2)+J_{n,2n}\right],\quad n=\pm1,\\
 P_c&=2(M_1+M_{-1})-S_3+4\zeta(3),\\
 P_f&=-J_{1,0}-J_{-1,0}+S_3+2\zeta(3).
\end{align}
Evaluating the $q\to0$ end of the same primitives gives
\begin{equation}
 K_a=3\zeta(3)+\operatorname{Li}_3(a^{-2})+\frac{P_A}{2},\quad
 K_c=-2S_3-P_c,\quad K_f=-4\zeta(3)-2S_3+P_f.
\end{equation}
The imaginary terms cancel in each final constant. The shuffle
identity gives $K_b=-K_c-K_f$.

For the outer letter $-4$, the only additional dilogarithmic value is
\begin{equation}
 \operatorname{Im}\operatorname{Li}_2(2a)
          =\frac56\mathcal C+\frac\pi2\ln2.
\end{equation}
One way to establish it is to use
$D(z)=\operatorname{Im}\operatorname{Li}_2(z)+\arg(1-z)\ln|z|$.
Its differential is
$\dd D=\ln|z|\,\dd\arg(1-z)-\ln|1-z|\,\dd\arg z$.
Integration gives one half of the sum of $\operatorname{Cl}_2$ at
twice the three angles of the triangle with vertices $0,1,z$.
For $z=2a$ those angles are $\pi/3,\pi/2,\pi/6$, and
$\arg(1-2a)=-\pi/2$. This yields the displayed value with its
branch fixed. In the paired endpoint primitives all
$\operatorname{Li}_3(2a)$ and $\operatorname{Li}_3(2a^{-1})$ terms
cancel. The remaining pair contributes
\begin{equation}
 -4\theta\pi\ln2+8\theta\operatorname{Im}\operatorname{Li}_2(2a)
     =\frac{20}{3}\theta\mathcal C.
\end{equation}
Together with $J(2,1)=-7\zeta(3)/4$, this fixes $K_e$.
The value of $J(2,1)$ follows from trilogarithm inversion and
$\operatorname{Li}_3(1/2)=7\zeta(3)/8-\pi^2\ln2/12+\ln^32/6$.

\paragraph{The remaining constant from the soft master.} The exact
energy equation \eqref{eq:comp-ck-I2-energy} determines $K_d$ once
the finite high-energy constant of $I_2$ is known. It can be extracted
without root-valued endpoint functions. At large $X$, use
\begin{equation}
 I_2^{\rm app}(X)=\int_0^1\frac{a}{Ud}
       \left[\frac12\ln X+\ln2+\ln U-\frac12\ln(1-U^2)\right]\dd U,
 \qquad I_2-I_2^{\rm app}=o(1).
\end{equation}
In this expression $a=\operatorname{artanh}(\beta U)$ and
$d=1-\beta^2U^2$ keep their exact $X$ dependence. To justify the
error estimate, split the interval at $U=X^{-1/4}$. The lower
pieces of both integrals are $O(X^{-1/4}\ln X)$. In the upper piece,
use the exact identity
\begin{equation}
 f=\ln\frac{(k+U)\sqrt s}{\sqrt{1-U^2}},\qquad
 \frac{k}{U}=\sqrt{1+\frac{1-U^2}{sU^2}}.
\end{equation}
Expanding these positive quantities bounds the integrated difference
by $O(X^{-3/4}\ln^2X)+O(X^{-1}\ln^2X)$. Thus the approximation
has the same finite constant. No approximation has been made to the
endpoint-regulating factor $d$.

Partial fractions and the fixed-letter differential reduction give
\begin{equation}
\begin{aligned}
 I_2^{\rm app}={}&\frac\ell4G(0,-1)-\frac\ell8G(-1,-1)
 +\frac12G(-2,0,-1)+\frac12G(0,-2,-1)\\
 &-\frac14G(-2,-1,-1)-\frac14G(0,-1,-1)
 -\frac14G(-1,-2,-1)\\
 &-\frac14G(-1,0,-1)+\frac14G(-1,-1,-1).
\end{aligned}
\end{equation}
Every endpoint is $X$. The few additional rational-letter constants
can be evaluated by elementary parameter differentiation. Let
$k_{012}(b)$ be the finite large-$X$ constant of $G(0,-1,-b;X)$.
Differentiating its defining integrals with respect to $b$ gives
\begin{equation}
 \partial_bG(0,-1,-b;X)
 =-\frac{\operatorname{Li}_2(-X)}{b(b-1)}
       +\frac{\operatorname{Li}_2(-X/b)}{b-1},
\end{equation}
so dilogarithm inversion yields
\begin{equation}
 k_{012}'(b)=-\frac{\ln^2b}{2(b-1)}-\frac{\zeta(2)}b.
\end{equation}
Since $k_{012}(1)=\zeta(3)$ and
$\int_1^2\ln^2b\,\dd b/(b-1)=\zeta(3)/4$, one obtains
$k_{012}(2)=7\zeta(3)/8-\zeta(2)\ln2$. The latter definite
integral follows from the trilogarithm value at one half stated above.
Explicitly it is
$2\zeta(3)-2\operatorname{Li}_3(1/2)-\pi^2\ln2/6+\ln^32/3$.
For $G(-2,-1,-1;X)$, shifting the integration variable by one and
then inverting it gives the constant
\begin{equation}
 -\frac12\int_1^\infty\frac{\ln^2v}{v(1+v)}\,\dd v
 =-\frac12\int_0^1\frac{\ln^2t}{1+t}\,\dd t
 =-\frac34\zeta(3).
\end{equation}
The last equality follows by integrating the alternating geometric
series. The shuffle product $G(-1)G(-2,-1)$ then fixes the constant
of $G(-1,-2,-1)$ to $3\zeta(3)/2$. Finally the shuffle identity
fixes the sum of the constants of $G(-2,0,-1)$ and $G(0,-2,-1)$
to $-7\zeta(3)/8$. Therefore
\begin{equation}
 I_2=\frac{5\ell^3}{48}+\frac{\pi^2\ell}{16}
                         -\frac38\zeta(3)+o(1).
\end{equation}
Only this constant is extracted with the approximation. The sharper
inverse-power remainder of the cross section follows from its exact
GPL kernels, rather than from the approximation error bound.

Collecting the endpoint evaluations gives
\begin{align}
 K_a&=\frac59\zeta(3)-\frac{4\pi}{9}\mathcal C,&
 K_b&=\frac23\zeta(3),\\
 K_c&=\frac23\zeta(3)-\frac{4\pi}{3}\mathcal C,&
 K_d&=\frac1{12}\zeta(3)+\frac\pi3\mathcal C,\\
 K_e&=\frac{16}{9}\zeta(3)-\frac{8\pi}{9}\mathcal C,&
 K_f&=-\frac43\zeta(3)+\frac{4\pi}{3}\mathcal C.
 \label{eq:comp-high-root-constants}
\end{align}
For $K_d$, the constant term of the exact $I_2$ identity is
$3(K_f+2K_a-K_e)/8+(K_d+K_c)/2=-3\zeta(3)/8$.
